% =====================================================================
% Theorem D(K) and Lemma K — single-cusp defect bounds for Bianchi groups
% Self-contained analysis paper: explicit, Arb-computable constants.
% NO gluing, NO N𝔭, NO reference cell, NO two-cusp coupling in main theorems.
% Revision after referee report 2026-07-12: full C_K proof, no Luke lower
% bound, consistent C_tr, metric comparison, non-circular η0, named C1/C2.
% =====================================================================
\documentclass[11pt]{amsart}
\usepackage{amsmath,amssymb,amsthm,mathtools,booktabs,hyperref,enumitem}
\usepackage{microtype}

\newtheorem{theorem}{Theorem}[section]
\newtheorem{lemma}[theorem]{Lemma}
\newtheorem{proposition}[theorem]{Proposition}
\newtheorem{corollary}[theorem]{Corollary}
\theoremstyle{definition}
\newtheorem{definition}[theorem]{Definition}
\newtheorem{assumption}[theorem]{Assumption}
\newtheorem{remark}[theorem]{Remark}

\newcommand{\HH}{\mathbb{H}^3}
\newcommand{\RR}{\mathbb{R}}
\newcommand{\CC}{\mathbb{C}}
\newcommand{\ZZ}{\mathbb{Z}}
\newcommand{\QQ}{\mathbb{Q}}
\newcommand{\NN}{\mathbb{N}}
\newcommand{\PSL}{\mathrm{PSL}}
\newcommand{\SL}{\mathrm{SL}}
\newcommand{\vol}{\mathrm{vol}}
\newcommand{\tr}{\mathrm{tr}}
\newcommand{\diam}{\mathrm{diam}}
\newcommand{\OK}{\mathcal{O}_K}
\newcommand{\NK}{N_{K/\QQ}}
\newcommand{\abs}[1]{\lvert#1\rvert}
\newcommand{\norm}[1]{\lVert#1\rVert}
\newcommand{\inner}[2]{\langle#1,#2\rangle}
\newcommand{\eps}{\varepsilon}
\newcommand{\tht}{\theta}

\title[Lemma K and Theorem D$(K)$]{Uniform $K$-Bessel tails and single-cusp
defect bounds for Bianchi groups\\
\large Explicit Arb-computable constants under a Hecke growth hypothesis}
\author{}
\date{Draft, 2026-07-12 (revision)}

\begin{document}
\begin{abstract}
We prove two foundational analytic statements for the Laplace spectrum of
Bianchi orbifolds $\Gamma\backslash\HH$, $\Gamma=\PSL_2(\OK)$, at class number
one ($K=\QQ(i)$ or $\QQ(\sqrt{-3})$).  Both theorems are stated under a
single, explicit Hecke growth hypothesis (Assumption~H), and every constant
is given by a closed formula that is enclosable in ball arithmetic (Arb /
python-flint).

\begin{enumerate}[leftmargin=*,itemsep=2pt]
\item \textbf{Lemma~K.}  An elementary pointwise \emph{upper} bound for
$|K_{ir}(y)|$ (proved from the integral representation, with no
$r$-independent lower bound claimed), combined with Assumption~H, yields an
explicit tail bound on the Fourier series of any cuspidal Hecke eigenform in
the cusp collar.  The majorant is Arb-computable and is validated against
truncated analytic sums at truncations $M=100,200,400$.
\item \textbf{Theorem~D$(K)$.}  A $K$-finite truncated Fourier expansion
with small automorphy defect and small $L^2$ residual is within
explicit distance of a true $L^2$ cusp form, with eigenvalue error
controlled by constants $C_1$ and $C_2$ assembled as named products of
trace, metric-comparison, Payne--Weinberger, Sobolev, and Lemma~K
contributions.  The threshold $\eta_0$ is defined from these a~priori
constants only (no circularity).
\end{enumerate}

\noindent\textbf{Known mathematics used:} modified-Bessel integral
representation and monotonicity in order, Rankin--Selberg / Kim--Sarnak
coefficient bounds (as inputs to~H), Euclidean trace and Poincar\'e
inequalities on tetrahedra with explicit hyperbolic-metric comparison on
$\{y\ge y_{\min}\}$, spectral resolution for cofinite Kleinian groups
(EGM / Friedman), and the Child / Booker--Str\"ombergsson--Venkatesh
defect philosophy (constants tracked).

\noindent\textbf{Engineering:} Arb enclosure of incomplete-gamma majorants;
floating comparison diagnostics (clearly labelled non-certifying).

\noindent\textbf{Not claimed here:} two-cusp coupling, congruence level
$N\mathfrak{p}$, reference-cell gluing, certified counting, Weyl law as an
exact count, or a uniform spectral gap $\mu(N)$.  Those are later rungs of
the dual-certification architecture; once Theorem~D$(K)$ is fixed, the
two-cusp version is a block-matrix estimate built on the same constants.
\end{abstract}
\maketitle

\tableofcontents

% ---------------------------------------------------------------------
\section{Introduction}
\label{sec:intro}

The Selberg-type eigenvalue problem for Bianchi groups asks for lower
bounds on the first positive Laplace eigenvalue $\lambda_1$ of
$\Gamma\backslash\HH$, $\Gamma=\PSL_2(\OK)$.  The natural threshold is
$\lambda_1\ge 1$ (the bottom of the continuous spectrum).  The best
general unconditional bound is
\[
  \lambda_1 \ge \frac{975}{1024}
\]
of Blomer--Brumley~\cite{BB11}, obtained from the Kim--Sarnak exponent
$7/64$ via Jacquet--Langlands.  At level one for the class-number-one
fields $K=\QQ(i)$ and $K=\QQ(\omega)$, numerical work of Then~\cite{Then}
predicts a first cuspidal eigenvalue with spectral parameter
\[
  r_1 \approx 6.62212,\qquad \lambda_1 = r_1^2+1 \approx 44.85,
\]
far above the threshold~$1$.  Making such values rigorous requires two
complementary ingredients: \emph{exclusion} of spectrum in windows that
contain no eigenvalues, and \emph{existence} of spectrum near a putative
eigenvalue.  The dual-certification program of the companion
roadmap~\cite{DualRoadmap} combines verified finite elements (exclusion /
counting) with interval Hejhal methods (existence).  Both sides need
explicit Fourier-tail and defect constants.  Those constants are the
subject of the present paper.

\paragraph{What this paper does.}
We isolate the single-cusp analytic core and prove it with fully
explicit constants.  No mesh, no gluing, and no congruence level appear
in the main theorems.  The two results are:

\begin{itemize}[leftmargin=*]
\item \textbf{Lemma~K} (Section~\ref{sec:lemmaK}): uniform $K_{ir}$ tails
under Hecke control of Fourier coefficients, from a fully elementary
pointwise \emph{upper} bound.
\item \textbf{Theorem~D$(K)$} (Section~\ref{sec:DK}): a defect bound
converting a nearly automorphic truncated expansion into a nearby true
cusp form, with every factor in $C_1,C_2$ named.
\end{itemize}

\paragraph{What this paper does not do.}
We do \emph{not} treat two-cusp coupling, the block matrix
$\bigl(\begin{smallmatrix}V_\infty&-S\\-S^*&V_0\end{smallmatrix}\bigr)$,
congruence covers, reference-cell tilings, certified counting, or Weyl
asymptotics used as exact counts.  An optional appendix records a
heuristic for $\mu(N)\ge c/\log N$; it is not used in any theorem.

\paragraph{Engineering target.}
Section~\ref{sec:numeric} records Then's value $r_1\approx 6.62212$ as an
\emph{engineering target} interval $[6.62211,6.62213]$.  We do
\emph{not} claim this as a theorem of the present paper: converting the
target into a certified eigenvalue requires a concrete defect $\eta$
small enough that $C_1\eta+C_2 e^{-2\pi Y_0}$ fits inside the desired
window, which is future computational work built on Theorem~D$(K)$.

% ---------------------------------------------------------------------
\section{Notation}
\label{sec:notation}

Let $K=\QQ(i)$ or $K=\QQ(\sqrt{-3})$, with ring of integers
$\OK=\ZZ[i]$ or $\OK=\ZZ[\omega]$, $\omega=e^{2\pi i/3}$.  Write
$\Gamma=\PSL_2(\OK)$ and let $\HH=\CC\times(0,\infty)$ be hyperbolic
$3$-space with metric and volume form
\begin{equation}
\label{eq:hypmetric}
  ds^2 = \frac{\abs{dz}^2+dy^2}{y^2},\qquad
  dV = \frac{dx_1\,dx_2\,dy}{y^3},\quad z=x_1+ix_2.
\end{equation}
The group $\Gamma$ has a single cusp (class number one), represented by
$\infty$.  A standard fundamental domain $F$ (Humbert for $\ZZ[i]$,
EGM / Swan for $\ZZ[\omega]$; cf.\ Elstrodt--Grunewald--Mennicke~\cite{EGM})
has cusp cross-section the flat torus (orbifold)
$T=\CC/\OK$ of Euclidean area
\begin{equation}
\label{eq:areaT}
  \abs{T}
  =
  \begin{cases}
    1/2, & \OK=\ZZ[i],\\
    \sqrt{3}/6, & \OK=\ZZ[\omega].
  \end{cases}
\end{equation}
(The $\ZZ[\omega]$ value is $\mathrm{covol}(\ZZ[\omega])/3=\sqrt{3}/6$,
matching the EGM planar section $P_3$.)

Fix truncation heights $Y\ge Y_0\ge 1/2$.  The truncated core and collar
are
\[
  K_Y := F\cap\{y\le Y\},\qquad
  C_Y := T\times[Y,\infty)
  \quad\bigl(\text{identified with }F\cap\{y\ge Y\}\bigr).
\]
Write $F_Y:=F\cap\{y\le Y\}$ for the closed truncated domain.  On $F_Y$
one has $y\ge y_{\min}(K)>0$; we use the conservative explicit floors
\begin{equation}
\label{eq:ymin}
  y_{\min}(\ZZ[i])
  :=
  \tfrac12,
  \qquad
  y_{\min}(\ZZ[\omega])
  :=
  \sqrt{\tfrac23},
\end{equation}
which lower-bound the standard Humbert / EGM floors
($y\ge\sqrt{1-\max\abs{z}^2}$ on the Ford spheres, with planar section
inside the unit cell).  The value $\tfrac12$ is strictly below
$1/\sqrt{2}$ (Humbert), and $\sqrt{2/3}$ is a standard conservative floor
for the Eisenstein--Picard domain; both may be sharpened without changing
the structure of the constants.

The positive Laplace--Beltrami operator on $\HH$ has continuous spectrum
$[1,\infty)$ on $\Gamma\backslash\HH$.  Cuspidal $L^2$ eigenfunctions
with eigenvalue $\lambda=r^2+1$, $r\in\RR$, admit a Fourier expansion at
the cusp $\infty$ of the form
\begin{equation}
\label{eq:fourier}
  f(z,y)
  =
  \sum_{0\neq\beta\in\OK}
  a_\beta\,
  y\,K_{ir}(2\pi\abs{\beta}y)\,
  e^{2\pi i\,\mathrm{Re}(\beta\overline{z})},
\end{equation}
with $K_{ir}$ the modified Bessel function of the second kind of pure
imaginary order (cf.\ EGM~\cite{EGM}, Friedman~\cite{Friedman}).  We write
$N(\beta)=\abs{\beta}^2$ for the field norm on principal ideals (so
$N(\beta)=\beta\overline{\beta}$ for $\beta\in\OK$), and
\[
  r_2(n) := \#\{\beta\in\OK:N(\beta)=n\}.
\]
For $\OK=\ZZ[i]$ one has the classical formula $r_2(n)=4(\chi_4*1)(n)$;
for $\OK=\ZZ[\omega]$ the corresponding representation numbers are
bounded by $6\,d(n)$.  In all estimates we use as theorem default the
divisor majorant
\begin{equation}
\label{eq:r2maj}
  r_2(n)\le 6\,d(n)\le 6n
  \qquad(n\ge 1),
\end{equation}
valid for both rings; the cruder $6n$ alone is retained as an optional
looser bound (same power of $n$ after absorption).

% ---------------------------------------------------------------------
\section{Assumption H}
\label{sec:H}

\begin{assumption}[H$(K,\tht)$]
\label{ass:H}
There exists a function $C_H:(0,\infty)\to(0,\infty)$ such that for every
$L^2$ cuspidal Hecke eigenform on $\Gamma\backslash\HH$ with Fourier
coefficients $a_\beta$ as in~\eqref{eq:fourier}, and for every
$\eps>0$,
\begin{equation}
\label{eq:H}
  \abs{a_\beta}
  \le
  C_H(\eps)\,
  N(\beta)^{\tht+\eps}
  \qquad\text{for all }\beta\in\OK\setminus\{0\}.
\end{equation}
\end{assumption}

\paragraph{Admissible exponents.}
\begin{itemize}[leftmargin=*]
\item $\tht=\tfrac12$ is available unconditionally from Rankin--Selberg
theory (cf.\ the $\mathrm{GL}_2$ literature over imaginary quadratic
fields; the constant $C_H(\eps)$ depends on the form's $L^2$-normalization
and on $\eps$).
\item $\tht=\tfrac7{64}$ is the Kim--Sarnak exponent~\cite{KS}, available
under the corresponding transfer to $\mathrm{GL}_2/K$.
\item $\tht=0$ is the Ramanujan conjecture at finite places (open in
general; recorded for comparison).
\end{itemize}

\begin{remark}
\label{rem:Hecke-entry}
Every constant appearing in Lemma~K and Theorem~D$(K)$ depends on
$(K,\tht)$ only through $C_H(\eps)$, $\tht$, and the geometric data
$(\abs{T},Y,Y_0,y_{\min})$.  There is no hidden dependence on a particular
eigenform beyond $C_H$.  In numerical work one may set $C_H=1$ after
absorbing normalization into the putative expansion, provided the
resulting $\eta$ is interpreted consistently.
\end{remark}

\paragraph{Standing convention.}
Every theorem below is understood to be stated under Assumption~H$(K,\tht)$
for a fixed admissible $\tht$ and a fixed choice of $C_H(\eps)$.  We write
$C_H$ for $C_H(\eps)$ once $\eps\ge 0$ has been fixed (the case $\eps=0$
is permitted whenever a polynomial bound without $\eps$ is known).

% ---------------------------------------------------------------------
\section{Assumptions A and S}
\label{sec:AS}

\begin{assumption}[A --- analyticity / unique continuation]
\label{ass:A}
Cuspidal $L^2$ eigenfunctions of $\Delta$ on $\Gamma\backslash\HH$ are
real-analytic on the compact core, and unique continuation holds across
$\overline{K_Y}$ (Morrey--Nirenberg / Aronszajn).  Consequently a nearly
automorphic putative form that is $L^2$-close to a true eigenfunction on
an open set of $F_Y$ coincides with it up to the defect controlled by
Theorem~\ref{thm:DK}.
\end{assumption}

\begin{assumption}[S --- spectral decomposition]
\label{ass:S}
The spectral resolution of $\Delta$ on $L^2(\Gamma\backslash\HH)$ for the
class-number-one Bianchi groups under consideration (trivial character)
is the standard discrete-plus-continuous decomposition of
Elstrodt--Grunewald--Mennicke~\cite{EGM} / Friedman~\cite{Friedman}
(Thm.~3.8.1 and related): the cuspidal subspace admits an orthonormal
basis of Maass cusp forms, and the continuous spectrum is spanned by
Eisenstein series (Lax--Phillips~\cite{LP}).
\end{assumption}

\paragraph{Standing convention (continued).}
Theorem~\ref{thm:DK} is stated under Assumptions~\ref{ass:H}, \ref{ass:A},
and~\ref{ass:S} simultaneously.  Lemma~K uses only Assumption~\ref{ass:H}.

% ---------------------------------------------------------------------
\section{Lemma K: uniform $K$-Bessel tails}
\label{sec:lemmaK}

\subsection{Pointwise upper bound (elementary)}

\begin{lemma}[Pointwise $K$-Bessel upper bound]
\label{lem:Kpoint}
For every $r\ge 0$ and every $y>0$,
\begin{equation}
\label{eq:Kpoint}
  \abs{K_{ir}(y)}
  \le
  \sqrt{\frac{\pi}{2y}}\,e^{-y}\,C_K(r),
\end{equation}
where one may take the sharp constant
\begin{equation}
\label{eq:CK}
  C_K(r)
  :=
  1
\end{equation}
independent of $r$.  In particular the classical (looser) choice
\begin{equation}
\label{eq:CKclass}
  C_K^{\mathrm{cl}}(r)
  :=
  \exp\Bigl(\frac{\pi r}{2}\Bigr)
  \ge 1
\end{equation}
also satisfies~\eqref{eq:Kpoint}; we record it only for comparison with the
automorphic literature.  All theorems below use the sharp value~\eqref{eq:CK}.

Moreover, for every $y\ge 1$ the first asymptotic correction for the
comparison function $K_0$ yields the optional \emph{upper} refinement
\begin{equation}
\label{eq:Luke}
  \abs{K_{ir}(y)}
  \le
  \sqrt{\frac{\pi}{2y}}\,e^{-y}\,
  \Bigl(1+\frac{1}{8y}\Bigr),
\end{equation}
so that $C_K^{\mathrm{poly}}:=1+(1/8)$ majorizes $\abs{K_{ir}(y)}$
uniformly for all $y\ge 1$ under the $K_0$-comparison of the proof
(alternatively, the $r$-dependent envelope
$C_K^{\mathrm{cl}}(r)\,(1+(1+r^2)/(4y))$ is a valid but much looser
upper bound).  We do \emph{not} claim any $r$-independent
\emph{lower} bound of the form $\sqrt{\pi/(2y)}\,e^{-y}\le\abs{K_{ir}(y)}$;
such a bound fails for large $r$ at fixed $y$.  Only upper bounds enter
the tail estimate.
\end{lemma}

\begin{proof}
We use the standard integral representation
(DLMF~\cite{DLMF}~\S10.32.9; Watson~\cite{Watson}~\S6.22), valid for all
$y>0$ and all real $r$:
\begin{equation}
\label{eq:Kint}
  K_{ir}(y)
  =
  \int_0^\infty e^{-y\cosh t}\,\cos(rt)\,dt.
\end{equation}
(The formula is the specialization $\nu=ir$ of
$K_\nu(y)=\int_0^\infty e^{-y\cosh t}\cosh(\nu t)\,dt$, using
$\cosh(irt)=\cos(rt)$; absolute convergence for $y>0$ is elementary.)

\paragraph{Step A (oscillation killed).}
Since $\abs{\cos(rt)}\le 1$ for all real $r,t$,
\begin{equation}
\label{eq:KA}
  \abs{K_{ir}(y)}
  \le
  \int_0^\infty e^{-y\cosh t}\,dt
  =
  K_0(y).
  \tag{A}
\end{equation}

\paragraph{Step B (monotonicity in order).}
For fixed $y>0$ the map $\nu\mapsto K_\nu(y)$ is even and strictly
increasing on $[0,\infty)$ (DLMF~\cite{DLMF}~\S10.37; alternatively, from
the integral representation
$K_\nu(y)=\int_0^\infty e^{-y\cosh t}\cosh(\nu t)\,dt$ and the fact that
$t\mapsto\cosh(\nu t)$ is increasing in $\nu\ge 0$ pointwise).  In
particular, with $\nu=\tfrac12$,
\begin{equation}
\label{eq:KB}
  K_0(y)
  \le
  K_{1/2}(y).
  \tag{B}
\end{equation}

\paragraph{Step C (closed form for half-order).}
The half-order modified Bessel function is elementary
(DLMF~\cite{DLMF}~\S10.39.2; Watson~\cite{Watson}~\S3.71):
\begin{equation}
\label{eq:KC}
  K_{1/2}(y)
  =
  \sqrt{\frac{\pi}{2y}}\,e^{-y}.
  \tag{C}
\end{equation}

\paragraph{Conclusion.}
Combining (A)--(C),
\[
  \abs{K_{ir}(y)}
  \le
  K_0(y)
  \le
  K_{1/2}(y)
  =
  \sqrt{\frac{\pi}{2y}}\,e^{-y},
\]
which is~\eqref{eq:Kpoint} with $C_K(r)=1$.  A fortiori the same holds
with $C_K^{\mathrm{cl}}(r)=e^{\pi r/2}\ge 1$.

\paragraph{Optional poly refinement for $y\ge 1$.}
The asymptotic expansion
$K_0(y)=\sqrt{\pi/(2y)}\,e^{-y}\,(1-1/(8y)+O(y^{-2}))$
and the classical inequality
$0<K_0(y)\le \sqrt{\pi/(2y)}\,e^{-y}\,(1+1/(8y))$
for $y\ge 1$ (cf.\ Luke~\cite{Luke}, or the remainder control in
DLMF~\cite{DLMF}~\S10.40) give~\eqref{eq:Luke} via Step~A.  This is an
upper refinement only; it is not used in the main tail algebra below.
\end{proof}

\begin{remark}
Both $C_K(r)\equiv 1$ and $C_K^{\mathrm{cl}}(r)$ are elementary and
Arb-enclosable.  The implementation \texttt{lemma\_K.py} exposes
\texttt{C\_K(r)} returning~$1$ by default, with
\texttt{classical=True} for~\eqref{eq:CKclass} and optional
\texttt{poly=True} for the $y\ge 1$ envelope.
\end{remark}

\subsection{Tail bound under Assumption H}

Define, for $M\ge 1$, $r\ge 0$, $Y_0>0$,
\begin{equation}
\label{eq:SMinf}
  S_{M,\infty}(r,Y_0)
  :=
  \sum_{N(\beta)>M}
  \abs{a_\beta}^2\,
  \abs{K_{ir}(2\pi\abs{\beta}Y_0)}^2.
\end{equation}

\begin{theorem}[Lemma K --- tail bound]
\label{thm:lemmaK}
Under Assumption~\ref{ass:H}, for every $M\ge 1$, $Y_0>0$, $r\ge 0$, and
every $\eps\ge 0$,
\begin{equation}
\label{eq:lemmaK}
\begin{aligned}
  S_{M,\infty}(r,Y_0)
  &\le
  \frac{C_H(\eps)^2\,C_K(r)^2}{4Y_0}
  \sum_{n>M}
  r_2(n)\,
  n^{2\tht+2\eps-\tfrac12}\,
  e^{-4\pi\sqrt{n}\,Y_0},
\end{aligned}
\end{equation}
with $C_K(r)=1$ as in~\eqref{eq:CK}.  Using the theorem-default multiplicity
majorant $r_2(n)\le 6\,d(n)$ and the further bound $d(n)\le n$,
\begin{equation}
\label{eq:lemmaKcrude}
  S_{M,\infty}(r,Y_0)
  \le
  \frac{3\,C_H(\eps)^2\,C_K(r)^2}{2Y_0}
  \sum_{n>M}
  n^{2\tht+2\eps+\tfrac12}\,
  e^{-4\pi\sqrt{n}\,Y_0}.
\end{equation}
(The same display holds with the optional looser $r_2(n)\le 6n$ used
directly.)  The series on the right of~\eqref{eq:lemmaKcrude} is bounded
by the incomplete-gamma comparison of Lemma~\ref{lem:integral} below,
yielding an explicit Arb ball
\[
  \eps_{\mathrm{tail}}(M,Y_0,r,\tht,\eps;C_H)
  \supset
  S_{M,\infty}(r,Y_0).
\]
\end{theorem}

\begin{proof}
By Lemma~\ref{lem:Kpoint} with $y=2\pi\abs{\beta}Y_0=2\pi\sqrt{n}\,Y_0$,
\begin{equation}
\label{eq:Ksq}
  \abs{K_{ir}(y)}^2
  \le
  \frac{\pi}{2y}\,e^{-2y}\,C_K(r)^2
  =
  \frac{1}{4\sqrt{n}\,Y_0}\,
  e^{-4\pi\sqrt{n}\,Y_0}\,
  C_K(r)^2.
\end{equation}
Assumption~H gives $\abs{a_\beta}^2\le C_H(\eps)^2\,n^{2\tht+2\eps}$.
Multiplying and summing over $\beta$ with $N(\beta)=n>M$, then over $n$,
produces~\eqref{eq:lemmaK}.  The bound $r_2(n)\le 6n$ upgrades the
power of $n$ by~$1$ and the prefactor by~$6$, giving~\eqref{eq:lemmaKcrude}
since $6/(4Y_0)=3/(2Y_0)$.
\end{proof}

\begin{lemma}[Integral comparison]
\label{lem:integral}
Let $c>0$, $\alpha\in\RR$, and $f(x)=x^\alpha e^{-c\sqrt{x}}$ for $x>0$.
The logarithmic derivative
$f'(x)/f(x)=\alpha/x - c/(2\sqrt{x})$ is negative for all
$x>(2\alpha/c)^2$ when $\alpha>0$, and for all $x\ge 1$ when
$\alpha\le 0$ and $c>0$.  Hence if
\begin{equation}
\label{eq:Nmono}
  N_0
  \ge
  N_{\mathrm{mono}}
  :=
  \begin{cases}
    \bigl\lceil\bigl(2\alpha/c\bigr)^2\bigr\rceil+1, & \alpha>0,\\
    1, & \alpha\le 0,
  \end{cases}
\end{equation}
then $f$ is positive and decreasing on $[N_0,\infty)$, and
\begin{equation}
\label{eq:intcomp}
  \sum_{n\ge N_0} f(n)
  \le
  f(N_0)
  +
  \int_{N_0}^\infty f(x)\,dx.
\end{equation}
The integral equals
\begin{equation}
\label{eq:intexact}
  \int_{N_0}^\infty x^\alpha e^{-c\sqrt{x}}\,dx
  =
  2\,c^{-(2\alpha+2)}\,
  \Gamma\bigl(2\alpha+2,\,c\sqrt{N_0}\bigr),
\end{equation}
where $\Gamma(s,z)=\int_z^\infty t^{s-1}e^{-t}\,dt$ is the upper incomplete
gamma function (convergent at infinity whenever the original integral
converges, i.e.\ for all real $\alpha$ since the exponential dominates).
\end{lemma}

\begin{proof}
The sign of $f'$ is elementary calculus, giving the monotonicity
threshold~\eqref{eq:Nmono}.  The comparison~\eqref{eq:intcomp} is the
standard integral-test remainder for a positive decreasing summand.  The
identity~\eqref{eq:intexact} follows from the substitution $u=\sqrt{x}$,
$x=u^2$, $dx=2u\,du$, then $t=cu$.
\end{proof}

\begin{remark}[Hecke input]
\label{rem:hecke-entry-K}
Assumption~H enters \emph{only} through the factors $C_H(\eps)^2$ and the
exponent $2\tht+2\eps$ in~\eqref{eq:lemmaK}.  Replacing $\tht$ by a smaller
exponent (Kim--Sarnak, Ramanujan) improves the power of $n$ but does not
change the exponential decay $e^{-4\pi\sqrt{n}Y_0}$, which dominates for
any fixed $Y_0>0$ as $M\to\infty$.
\end{remark}

\subsection{Arb implementation and benchmarks}

The companion module \texttt{lemma\_K.py} implements:
\begin{itemize}[leftmargin=*]
\item \texttt{C\_K(r)} --- Arb enclosure of~\eqref{eq:CK} (default~$1$);
\item \texttt{lemma\_K\_tail(M, Y0, r, theta, C\_H=1.0, eps=0.0)} --- Arb
enclosure of the right-hand side of~\eqref{eq:lemmaKcrude} via
Lemma~\ref{lem:integral} (default multiplicity mode \texttt{div}:
$r_2\le 6d(n)\le 6n$);
\item \texttt{C1\_C2\_constants(...)} --- Arb/float evaluation of the
named constants of Section~\ref{sec:DK};
\item benchmarks at $M=100,200,400$ comparing the infinite-tail enclosure
to a long truncated majorant sum that uses the \emph{exact} $r_2$ for
$\ZZ[i]$ (analytic majorant only; no fake Fourier coefficients).
\end{itemize}

Representative output at the engineering parameters
$(Y_0,r,\tht)=(0.8,6.6,1/2)$, $C_H=1$, $\eps=0$, $C_K=1$ (python-flint,
certifying): the tail at $M=400$ is on the order of $10^{-82}$ (see
\texttt{constants\_DK.md}); the classical $C_K^{\mathrm{cl}}$ inflates
the same majorant by $C_K^{\mathrm{cl}}(6.6)^2\approx 1.01\cdot 10^{9}$.

% ---------------------------------------------------------------------
\section{Theorem D$(K)$: single-cusp defect bound}
\label{sec:DK}

\subsection{Putative forms and defects}

\begin{definition}[Truncated putative form]
\label{def:putative}
Fix $r\ge 0$, $M\ge 1$, $Y_0\ge 1/2$, and complex coefficients
$(a_\beta)_{0<\abs{\beta}\le\sqrt{M}}$ (indexed by norm $N(\beta)\le M$).
A \emph{truncated putative form} is the $K$-finite function defined on the
collar by
\begin{equation}
\label{eq:putative}
  \widetilde f(z,y)
  =
  \sum_{0<N(\beta)\le M}
  a_\beta\,
  y\,K_{ir}\bigl(2\pi\abs{\beta}y\bigr)\,
  e^{2\pi i\,\mathrm{Re}(\beta\overline{z})}
  \qquad(y\ge Y_0),
\end{equation}
and extended to the truncated core $K_Y$ ($Y\ge Y_0$) by a fixed $C^2$
partition of unity subordinate to a finite cover of $\overline{K_Y}$
compatible with the side-pairing of $F$.  (The particular partition
affects only the chart-overlap factor absorbed into $C_{\mathrm{tr}}^{\mathrm{cons}}$
below; one may take a standard Whitney-type cutoff.)
\end{definition}

Let $S\subset\Gamma$ be a finite set of side-pairing generators for the
EGM / Humbert fundamental domain $F$ (translations generating $\OK$, the
inversion $S:z\mapsto -1/z$, and unit rotations when present).

\begin{definition}[Defects]
\label{def:defects}
The \emph{automorphy defect} and \emph{$L^2$ residual} of $\widetilde f$
are
\begin{align}
  \delta_{\mathrm{aut}}
  &:=
  \max_{\gamma\in S}
  \sup_{z\in F_Y}
  \abs{\widetilde f(\gamma z)-\widetilde f(z)},
  \label{eq:deltaaut}\\
  \tau_{\mathrm{tail}}
  &:=
  \norm{(\Delta-\lambda)\widetilde f}_{L^2_{\mathrm{hyp}}(K_Y)},
  \qquad\lambda:=r^2+1.
  \label{eq:tautil}
\end{align}
Set $\eta:=\delta_{\mathrm{aut}}+\tau_{\mathrm{tail}}$.
\end{definition}

\subsection{Geometric inequalities with explicit constants}

We record the elementary inequalities used to convert boundary defects
into $L^2$ control.  All constants are fully explicit.  Trace and
Poincar\'e are Euclidean; hyperbolic weights are converted in
Lemma~\ref{lem:metric}.

\begin{lemma}[Trace inequality on a tetrahedron]
\label{lem:T}
Let $T\subset\RR^3$ be a tetrahedron, $F$ a face of $T$, and $P$ the
opposite vertex.  For every $w\in H^1(T)$,
\begin{equation}
\label{eq:trace}
  \int_F w^2\,d\sigma
  \le
  \frac{\abs{F}}{\abs{T}}
  \Bigl(
    \norm{w}_{L^2(T)}^2
    +
    \tfrac23 h_T\,\norm{w}_{L^2(T)}\,\norm{\nabla w}_{L^2(T)}
  \Bigr),
\end{equation}
where $h_T=\diam(T)$ and all norms on the right are Euclidean.  Applying
Young's inequality in the form
\begin{equation}
\label{eq:Young}
  \tfrac23 h_T\,\norm{w}\,\norm{\nabla w}
  \le
  \tfrac13\norm{w}^2
  +
  \tfrac13 h_T^2\norm{\nabla w}^2
\end{equation}
yields
\begin{equation}
\label{eq:trace2}
  \norm{w}_{L^2(F)}^2
  \le
  C_{\mathrm{tr}}(T,F)
  \Bigl(
    \norm{w}_{L^2(T)}^2
    +
    h_T^2\norm{\nabla w}_{L^2(T)}^2
  \Bigr),
\end{equation}
with the single consistent constant
\begin{equation}
\label{eq:Ctr}
  C_{\mathrm{tr}}(T,F)
  :=
  \frac43\,\frac{\abs{F}}{\abs{T}}.
\end{equation}
\end{lemma}

\begin{proof}
The identity
$\operatorname{div}\bigl((x-P)w^2\bigr)=3w^2+2w(x-P)\cdot\nabla w$
integrates over $T$.  On faces containing $P$ one has $(x-P)\cdot\nu=0$;
on $F$ one has $(x-P)\cdot\nu=3\abs{T}/\abs{F}$.  Cauchy--Schwarz and
$\abs{x-P}\le h_T$ give~\eqref{eq:trace}.  Substitute~\eqref{eq:Young}:
\begin{align*}
  \int_F w^2
  &\le
  \frac{\abs{F}}{\abs{T}}
  \Bigl(
    \norm{w}^2
    +
    \tfrac13\norm{w}^2
    +
    \tfrac13 h_T^2\norm{\nabla w}^2
  \Bigr)
  =
  \frac{\abs{F}}{\abs{T}}
  \Bigl(
    \tfrac43\norm{w}^2
    +
    \tfrac13 h_T^2\norm{\nabla w}^2
  \Bigr)\\
  &\le
  \frac43\frac{\abs{F}}{\abs{T}}
  \Bigl(
    \norm{w}^2
    +
    h_T^2\norm{\nabla w}^2
  \Bigr),
\end{align*}
which is~\eqref{eq:trace2}--\eqref{eq:Ctr}.
\end{proof}

\begin{lemma}[Payne--Weinberger / Bebendorf]
\label{lem:PW}
Let $\Omega\subset\RR^d$ be bounded, convex, and open.  The first positive
Neumann eigenvalue of the Euclidean Laplacian on $\Omega$ satisfies
\begin{equation}
\label{eq:PW}
  \lambda_1^N(\Omega)
  \ge
  \frac{\pi^2}{\diam(\Omega)^2}.
\end{equation}
Equivalently, for every $w\in H^1(\Omega)$ with $\int_\Omega w=0$,
\begin{equation}
\label{eq:Poincare}
  \norm{w}_{L^2(\Omega)}
  \le
  \frac{\diam(\Omega)}{\pi}\,
  \norm{\nabla w}_{L^2(\Omega)},
\end{equation}
where the gradient is Euclidean.
\end{lemma}

\begin{proof}
Payne--Weinberger~\cite{PW60}; corrected proof by Bebendorf~\cite{Beb03}.
Both statements are for the Euclidean Laplacian on a convex Euclidean
domain; hyperbolic weights are handled by Lemma~\ref{lem:metric}.
\end{proof}

\begin{lemma}[Sobolev / mean control on $F_Y$]
\label{lem:Sobolev}
Let $D_Y:=\diam_{\mathrm{euc}}(F_Y)$ be the Euclidean diameter of a
convex body containing $F_Y$ (e.g.\ the axis-aligned box
$[-1,1]\times[-1,1]\times[y_{\min},Y]$).  For every $w\in H^1(F_Y)$ write
$\overline{w}=\vol_{\mathrm{euc}}(F_Y)^{-1}\int_{F_Y}w$.  Then
\begin{equation}
\label{eq:Sob2}
  \norm{w-\overline{w}}_{L^2_{\mathrm{euc}}(F_Y)}
  \le
  \frac{D_Y}{\pi}\,
  \norm{\nabla_{\mathrm{euc}} w}_{L^2_{\mathrm{euc}}(F_Y)},
\end{equation}
and the mean satisfies
$\abs{\overline{w}}\le \vol_{\mathrm{euc}}(F_Y)^{-1/2}\norm{w}_{L^2_{\mathrm{euc}}}$.
We set
\begin{equation}
\label{eq:CPoincSob}
  C_{\mathrm{Poincar\'e}}
  :=
  \frac{D_Y}{\pi},
  \qquad
  C_{\mathrm{Sob}}
  :=
  1+C_{\mathrm{Poincar\'e}},
\end{equation}
so that
$\norm{w}_{L^2_{\mathrm{euc}}}
\le C_{\mathrm{Sob}}\,
\bigl(\norm{w}_{L^2_{\mathrm{euc}}}^2
+\norm{\nabla_{\mathrm{euc}}w}_{L^2_{\mathrm{euc}}}^2\bigr)^{1/2}$
up to the harmless volume normalization of the mean term (absorbed into
the conservative $D_Y$ below).

A conservative explicit choice for the Picard and Eisenstein--Picard
cores at $Y=1.25$ is
\begin{equation}
\label{eq:diam}
  D_Y \le 10,
\end{equation}
justified by embedding $F_Y$ in the box of side-lengths
$(1,1/2,Y-y_{\min})$ for $\ZZ[i]$ (Humbert: $x_1\in[-1/2,1/2]$,
$x_2\in[0,1/2]$, $y\in[y_{\min},Y]$) or the corresponding EGM prism for
$\ZZ[\omega]$, both of which have Euclidean diameter $<3$ at $Y=1.25$;
the bound $10$ leaves a wide safety margin for orbifold charts and
partition-of-unity overlaps.
\end{lemma}

\begin{lemma}[Hyperbolic vs Euclidean comparison on $F_Y$]
\label{lem:metric}
On $F_Y\subset\{y_{\min}\le y\le Y\}$, the hyperbolic and Euclidean
structures~\eqref{eq:hypmetric} satisfy, for every measurable $w$ and
every $w\in H^1(F_Y)$,
\begin{align}
  Y^{-3/2}\,\norm{w}_{L^2_{\mathrm{euc}}}
  &\le
  \norm{w}_{L^2_{\mathrm{hyp}}}
  \le
  y_{\min}^{-3/2}\,\norm{w}_{L^2_{\mathrm{euc}}},
  \label{eq:M-L2}\\
  Y^{-1/2}\,\norm{\nabla_{\mathrm{euc}}w}_{L^2_{\mathrm{euc}}}
  &\le
  \norm{\nabla_{\mathrm{hyp}}w}_{L^2_{\mathrm{hyp}}}
  \le
  y_{\min}^{-1/2}\,\norm{\nabla_{\mathrm{euc}}w}_{L^2_{\mathrm{euc}}}.
  \label{eq:M-grad}
\end{align}
(The gradient comparison uses
$\abs{\nabla_{\mathrm{hyp}}w}=y\abs{\nabla_{\mathrm{euc}}w}$ and
$dV_{\mathrm{hyp}}=y^{-3}\,dV_{\mathrm{euc}}$, so
$\abs{\nabla_{\mathrm{hyp}}w}^2\,dV_{\mathrm{hyp}}
=y^{-1}\abs{\nabla_{\mathrm{euc}}w}^2\,dV_{\mathrm{euc}}$.)
Consequently we record the explicit metric factors
\begin{equation}
\label{eq:Cmet}
\begin{aligned}
  C_{\mathrm{met},0}
  &:=
  y_{\min}^{-3/2},
  \qquad
  C_{\mathrm{met},0}^{-}
  :=
  Y^{3/2},\\
  C_{\mathrm{met},1}
  &:=
  y_{\min}^{-1/2},
  \qquad
  C_{\mathrm{met},1}^{-}
  :=
  Y^{1/2},\\
  C_{\mathrm{met}}
  &:=
  \max\bigl(C_{\mathrm{met},0},\,C_{\mathrm{met},1},\,
  C_{\mathrm{met},0}^{-},\,C_{\mathrm{met},1}^{-}\bigr).
\end{aligned}
\end{equation}
\end{lemma}

\begin{proof}
Immediate from $y\in[y_{\min},Y]$ and~\eqref{eq:hypmetric}.
\end{proof}

\begin{definition}[Explicit geometric constants]
\label{def:geo-const}
For use in Theorem~\ref{thm:DK} we set
\begin{align}
  C_{\mathrm{trace}}(K,Y)
  &:=
  \max_{T\subset\mathcal{T}_Y}
  \frac43\frac{\abs{F_T}}{\abs{T}},
  \label{eq:Ctrace}\\
  C_{\mathrm{Poincar\'e}}(K,Y)
  &:=
  \frac{D_Y}{\pi},
  \label{eq:CPoinc}\\
  C_{\mathrm{Sob}}(K,Y)
  &:=
  1+C_{\mathrm{Poincar\'e}}(K,Y),
  \label{eq:CSob}
\end{align}
where $\mathcal{T}_Y$ is any fixed tetrahedralization of a convex body
containing $F_Y$, and $F_T$ runs over faces of $T$ that meet
$\partial F_Y$.  If no mesh is fixed, the conservative majorant
\begin{equation}
\label{eq:Ctrace-cons}
  C_{\mathrm{trace}}^{\mathrm{cons}}
  :=
  \frac{4}{h_{\min}},
  \qquad
  h_{\min}:=0.25
\end{equation}
is admissible: for a right tetrahedron of altitude $h$ one has
$\abs{F}/\abs{T}=3/h$, so
$\tfrac43\cdot 3/h=4/h$, and $h_{\min}=0.25$ undercuts the vertical
thickness $Y-y_{\min}$ and the planar inradius of the Humbert / EGM
sections at $Y=1.25$.  Numerically
$C_{\mathrm{trace}}^{\mathrm{cons}}=16$.  This is a
\emph{conservative geometric majorant}, not a sharp face/volume ratio of
the orbifold fundamental domain.
\end{definition}

\subsection{Cusp modes via Lemma K and Lax--Phillips}

\begin{lemma}[Collar control]
\label{lem:collar}
Let $f$ be a true $L^2$ cusp form with eigenvalue $\lambda=r^2+1$ and
Fourier coefficients $a_\beta$.  Under Assumption~H, for every
$Y_0\ge 1/2$,
\begin{equation}
\label{eq:collar}
  \norm{f}_{L^2(C_{Y_0})}^2
  \le
  \abs{T}\,
  \frac{C_H(\eps)^2\,C_K(r)^2}{4Y_0}
  \sum_{n\ge 1}
  r_2(n)\,
  n^{2\tht+2\eps-\tfrac12}\,
  e^{-4\pi\sqrt{n}\,Y_0}
  \cdot
  \int_{Y_0}^\infty y^{-1}\,e^{-4\pi\sqrt{n}\,(y-Y_0)}\,dy.
\end{equation}
In particular the contribution of modes with $N(\beta)>M$ is at most
$\abs{T}$ times a constant multiple of the right-hand side
of~\eqref{eq:lemmaK}.  The zero mode of a cusp form vanishes identically
(cuspidality), so Lax--Phillips reduction~\cite{LP} contributes no
additional residual from the constant term.
\end{lemma}

\begin{proof}
Parseval on $T$ reduces $\norm{f}_{L^2(C_{Y_0})}^2$ to a sum over
$\beta\neq 0$ of $\int_{Y_0}^\infty\abs{a_\beta\,y K_{ir}(2\pi\abs{\beta}y)}^2 y^{-3}\,dy$.
The standard identity for the Whittaker/$K$-Bessel profile
(cf.\ EGM~\cite{EGM}~\S3) and Lemma~\ref{lem:Kpoint} bound each mode by
the corresponding term in~\eqref{eq:collar}.  Cuspidality kills the zero
mode; Lax--Phillips~\cite{LP} identifies the continuous-spectrum
contribution with Eisenstein series, which are excluded from the cuspidal
subspace under consideration.
\end{proof}

\subsection{Statement of Theorem D$(K)$}

\begin{theorem}[Theorem D$(K)$]
\label{thm:DK}
Let $K=\QQ(i)$ or $\QQ(\sqrt{-3})$, $\Gamma=\PSL_2(\OK)$, and fix
geometric parameters
\[
  Y\ge Y_0\ge \tfrac12,\qquad r\ge 0,\qquad M\ge 1
\]
(truncation heights, spectral parameter, and Fourier cutoff).  Under
Assumptions~\ref{ass:H}, \ref{ass:A}, and~\ref{ass:S}, there exist explicit
constants
\begin{align}
  C_1 &= C_1(K,Y,Y_0,r,\tht;\eps,C_H,M),
  \label{eq:C1name}\\
  C_2 &= C_2(K,Y),
  \label{eq:C2name}
\end{align}
given by the formulas~\eqref{eq:C1}--\eqref{eq:C2} below and enclosable in
Arb arithmetic, and a threshold $\eta_0=\eta_0(K,Y,r)>0$ defined
from~\eqref{eq:eta0}, such that the following holds.

Let $\widetilde f$ be a truncated putative form as in
Definition~\ref{def:putative}, with $\lambda=r^2+1$ and defects
$\eta=\delta_{\mathrm{aut}}+\tau_{\mathrm{tail}}$ as in
Definition~\ref{def:defects}.  Write $U_{\mathrm{cusp}}$ for the
orthogonal projection of the periodized corrected function $U$
(constructed in the proof) onto the cuspidal subspace of
$L^2(\Gamma\backslash\HH)$.  If $\eta\le\eta_0$ and
$\norm{\widetilde f}_{L^2_{\mathrm{hyp}}(F_Y)}=1$, then there exists a true
$L^2$ cusp form $f$ on $\Gamma\backslash\HH$ with eigenvalue
$\lambda_{\mathrm{true}}$ satisfying
\begin{align}
  \abs{\lambda_{\mathrm{true}}-\lambda}
  &\le
  C_1\,\eta
  +
  C_2\,e^{-2\pi Y_0}
  =
  C_1\,(\delta_{\mathrm{aut}}+\tau_{\mathrm{tail}})
  +
  C_2\,e^{-2\pi Y_0},
  \label{eq:lamerr}\\
  \norm{U_{\mathrm{cusp}}-f}_{L^2(\Gamma\backslash\HH)}
  &\le
  C_1\sqrt{\eta}.
  \label{eq:L2err}
\end{align}
In particular, any later interval-Hejhal computation that produces certified
enclosures of $(\delta_{\mathrm{aut}},\tau_{\mathrm{tail}})$ at fixed
$(Y,Y_0,r,M)$ may plug those values into~\eqref{eq:lamerr} (via the Arb
API \texttt{defect\_to\_lambda\_error} of \texttt{defect\_bound\_arb.py})
and obtain an enclosure of $\abs{\lambda_{\mathrm{true}}-\lambda}$; the
defects enter only through the linear combination $C_1(\delta+\tau)$ plus
the explicit collar term $C_2 e^{-2\pi Y_0}$.
(The difference between $\widetilde f$ and $U_{\mathrm{cusp}}$ on $F_Y$ is
controlled by $A_{\mathrm{bdry}}\delta_{\mathrm{aut}}$ plus the collar
cutoff; see Step~5.)  Moreover,
\begin{equation}
\label{eq:C1order}
  C_1
  =
  O_{K,Y_0}\!\Bigl(
    (1+r^{2})\,
    C_{\mathrm{met}}\,
    C_{\mathrm{trace}}\,
    C_{\mathrm{Sob}}
  \Bigr),
\end{equation}
with no hidden $N\mathfrak{p}$ or gluing dependence.  (The factor $1+r^{2}$
comes from $A_{\mathrm{ell}}=1+\lambda$ in Step~2; the $H^1$ corrector
constant of Step~1 is $\lambda$-free.  With the sharp $C_K=1$ there is no
$e^{\pi r/2}$ factor in $C_1$; the classical looser $C_K^{\mathrm{cl}}$
reintroduces it only inside $A_{\mathrm{cusp}}$.)
\end{theorem}

\subsection{Named intermediate constants and formulas for $C_1$, $C_2$}

Write $\lambda=r^2+1$.  All of the following are a~priori geometric /
analytic constants (no dependence on $\eta$ or on $C_1$ itself).

\begin{align}
  A_{\mathrm{met}}
  &:=
  C_{\mathrm{met}}
  =
  \max\bigl(y_{\min}^{-3/2},\,y_{\min}^{-1/2},\,Y^{3/2},\,Y^{1/2}\bigr),
  \label{eq:Amet}\\
  A_{\mathrm{bdry}}
  &:=
  C_{\mathrm{trace}}(K,Y)\,
  C_{\mathrm{Sob}}(K,Y)\,
  A_{\mathrm{met}},
  \label{eq:Abdry}\\
  A_{\mathrm{ell}}
  &:=
  1+\lambda,
  \label{eq:Aell}\\
  A_{\mathrm{res}}
  &:=
  A_{\mathrm{ell}}\,
  \bigl(
    1+
    C_{\mathrm{Poincar\'e}}(K,Y)\,
    A_{\mathrm{met}}
  \bigr),
  \label{eq:Ares}\\
  A_{\mathrm{cusp}}
  &:=
  \frac{C_H(\eps)\,C_K(r)}{2\sqrt{Y_0}}\,
  \Bigl(
    \sum_{n>M}
    6n\cdot
    n^{2\tht+2\eps-\tfrac12}\,
    e^{-4\pi\sqrt{n}\,Y_0}
  \Bigr)^{1/2},
  \label{eq:Acusp}\\
  A_{\mathrm{cut}}
  &:=
  2\,\abs{T}\,
  Y^{-1}\,
  \bigl(1+C_{\mathrm{Sob}}(K,Y)\bigr)\,
  A_{\mathrm{met}}.
  \label{eq:Acut}
\end{align}
Here $A_{\mathrm{cusp}}$ is enclosed by Lemma~\ref{lem:integral} (and is
numerically negligible once $M\ge 100$ at $Y_0=0.8$).  The series already
contains $C_K$ through~\eqref{eq:Acusp}; we do \emph{not} multiply by an
extra $(1+C_K)$ factor.

Define
\begin{align}
  C_1(K,Y,Y_0,r,\tht)
  &:=
  2\,
  A_{\mathrm{bdry}}\,
  A_{\mathrm{res}}\,
  \bigl(1+A_{\mathrm{cusp}}\bigr),
  \label{eq:C1}\\
  C_2(K,Y)
  &:=
  A_{\mathrm{cut}}.
  \label{eq:C2}
\end{align}
The prefactor $2$ in~\eqref{eq:C1} is derived in Step~4 (normalization
$\norm{U}\ge\tfrac12$ under $\eta\le\eta_0$).  The threshold $\eta_0$ is
defined from these a~priori constants by
\begin{equation}
\label{eq:eta0}
  \eta_0
  :=
  \min\Bigl(
    \tfrac{1}{4(1+A_{\mathrm{bdry}})^2},\,
    \tfrac{1}{4\,C_1^{2}}
  \Bigr).
\end{equation}
(The first term guarantees $\norm{e_{\mathrm{bdry}}}_{H^1}\le\tfrac14$ and
hence $\norm{U}\ge\tfrac12$; the second guarantees residual
$<\tfrac12\norm{U}$ for the spectral pigeonhole.  Both use only constants
already defined.)

\begin{remark}[Where each factor comes from]
\label{rem:factor-map}
\begin{itemize}[leftmargin=*]
\item $C_{\mathrm{trace}}$: Lemma~\ref{lem:T} with~\eqref{eq:Ctr}.
\item $C_{\mathrm{Poincar\'e}},C_{\mathrm{Sob}}$: Lemmas~\ref{lem:PW}--\ref{lem:Sobolev}.
\item $A_{\mathrm{met}}$: Lemma~\ref{lem:metric}.
\item $A_{\mathrm{bdry}}$: Step~1 (automorphy $\to H^1$ corrector).
\item $A_{\mathrm{ell}},A_{\mathrm{res}}$: Step~2 (residual / elliptic).
\item $A_{\mathrm{cusp}}$: Theorem~\ref{thm:lemmaK} + Lemma~\ref{lem:collar}
(Hecke enters \emph{only} here, via $C_H,\tht$).
\item $A_{\mathrm{cut}}=C_2$: Step~3 (collar cutoff error).
\item Factor $2$ in $C_1$: Step~4 ($\norm{U}\ge 1/2$).
\end{itemize}
\end{remark}

\subsection{Proof of Theorem~\ref{thm:DK}}

The argument follows the spectral-resolution philosophy of
Booker--Str\"ombergsson--Venkatesh~\cite[Prop.~3.1]{BSV} and
Child~\cite[Thm.~1.1]{Child}, adapted to $\HH$ with every constant tracked.

\paragraph{Step 1: automorphy defect $\to$ boundary $L^2$ $\to$ $H^1$ corrector.}
Fix $\gamma\in S$ and a face $F$ of $F_Y$ paired by $\gamma$.  The
pointwise bound $\abs{\widetilde f(\gamma z)-\widetilde f(z)}\le\delta_{\mathrm{aut}}$
gives the Euclidean face bound
\begin{equation}
\label{eq:step1a}
  \norm{\widetilde f\circ\gamma-\widetilde f}_{L^2_{\mathrm{euc}}(F)}
  \le
  \delta_{\mathrm{aut}}\,
  \abs{F}^{1/2}.
  \tag{1a}
\end{equation}
Apply Lemma~\ref{lem:T} in reverse on a collar tetrahedron $T_F\subset F_Y$
of altitude $\ge h_{\min}$ meeting $F$: there exists an extension
$e_F\in H^1(T_F)$ of the jump, supported in $T_F$, with
\begin{equation}
\label{eq:step1b}
  \norm{e_F}_{L^2_{\mathrm{euc}}(T_F)}^2
  +
  h_T^2\norm{\nabla_{\mathrm{euc}}e_F}_{L^2_{\mathrm{euc}}(T_F)}^2
  \le
  C_{\mathrm{tr}}(T_F,F)^{-1}\,
  \norm{\mathrm{jump}}_{L^2(F)}^2
  \cdot C_{\mathrm{tr}}(T_F,F)^{2}
  \le
  C_{\mathrm{trace}}\,
  \abs{F}\,
  \delta_{\mathrm{aut}}^2,
  \tag{1b}
\end{equation}
(the standard right-inverse of the trace map on a Lipschitz tetrahedron,
with norm $\le C_{\mathrm{tr}}$; cf.\ the construction by reflection across
$F$ and cutoff at height $h_{\min}$).  Summing over the finitely many
faces and applying a $C^2$ partition of unity subordinate to
$\{T_F\}$ (overlap multiplicity $\le N_{\mathrm{ov}}\le 4$ for a standard
tetrahedralization, absorbed into the conservative
$C_{\mathrm{trace}}^{\mathrm{cons}}=16$) produces
$e_{\mathrm{bdry}}\in H^1(F_Y)$ supported in a tubular neighbourhood of
$\partial F_Y$ with
\begin{equation}
\label{eq:step1c}
  \norm{e_{\mathrm{bdry}}}_{H^1_{\mathrm{euc}}(F_Y)}
  \le
  C_{\mathrm{trace}}\,
  C_{\mathrm{Sob}}\,
  \delta_{\mathrm{aut}}.
  \tag{1c}
\end{equation}
Convert to hyperbolic norms by Lemma~\ref{lem:metric}:
\begin{equation}
\label{eq:step1}
  \norm{e_{\mathrm{bdry}}}_{H^1_{\mathrm{hyp}}(F_Y)}
  \le
  A_{\mathrm{met}}\,
  C_{\mathrm{trace}}\,
  C_{\mathrm{Sob}}\,
  \delta_{\mathrm{aut}}
  =
  A_{\mathrm{bdry}}\,\delta_{\mathrm{aut}}.
  \tag{1}
\end{equation}
In particular $A_{\mathrm{bdry}}$ is \emph{independent of $\lambda$}: it is
purely the product of the Euclidean trace/Sobolev constants with the
metric comparison on $\{y\ge y_{\min}\}$.  (Earlier drafts inserted an
extra $(1+\sqrt{\lambda})$ into $A_{\mathrm{bdry}}$ ``for elliptic
control''; that factor is already accounted for once in
$A_{\mathrm{ell}}=1+\lambda$ of Step~2, so keeping it in both places was
a double count; see Remark~\ref{rem:no-double-sqrt}.)

\paragraph{Step 2: residual and approximate eigenfunction.}
Set $u:=\widetilde f-e_{\mathrm{bdry}}$ on $F_Y$.  By construction the
jumps of $u$ across side-pairings vanish in the $H^{1/2}$ trace sense, so
$u$ descends to a single-valued $H^1$ function on the truncated orbifold
$\Gamma\backslash F_Y$.  The residual expands as
\begin{equation}
\label{eq:step2a}
  (\Delta-\lambda)u
  =
  (\Delta-\lambda)\widetilde f
  -
  (\Delta-\lambda)e_{\mathrm{bdry}},
  \tag{2a}
\end{equation}
so
\begin{equation}
\label{eq:step2b}
  \norm{(\Delta-\lambda)u}_{L^2_{\mathrm{hyp}}(K_Y)}
  \le
  \tau_{\mathrm{tail}}
  +
  \norm{(\Delta-\lambda)e_{\mathrm{bdry}}}_{L^2_{\mathrm{hyp}}}.
  \tag{2b}
\end{equation}
On the compact core $K_Y\subset\{y\ge y_{\min}\}$, the hyperbolic Laplacian
in coordinates satisfies the pointwise bound
\begin{equation}
\label{eq:step2c}
  \abs{\Delta_{\mathrm{hyp}} e}
  \le
  C_{\Delta}\,
  \bigl(
    \abs{\nabla_{\mathrm{euc}}^2 e}
    +
    y_{\min}^{-1}\abs{\nabla_{\mathrm{euc}} e}
  \bigr)
  \tag{2c}
\end{equation}
with $C_{\Delta}\le 3$ (from
$\Delta_{\mathrm{hyp}}=-y^2(\partial_{x_1}^2+\partial_{x_2}^2+\partial_y^2)+y\partial_y$).
For the Whitney corrector $e_{\mathrm{bdry}}$ built from face jumps of size
$\delta_{\mathrm{aut}}$ on tetrahedra of altitude $h_{\min}$, standard
scaling gives
\begin{equation}
\label{eq:step2d}
  \norm{e_{\mathrm{bdry}}}_{H^2_{\mathrm{euc}}}
  \le
  h_{\min}^{-1}\,
  \norm{e_{\mathrm{bdry}}}_{H^1_{\mathrm{euc}}}
  \le
  4\,
  C_{\mathrm{trace}}\,
  C_{\mathrm{Sob}}\,
  \delta_{\mathrm{aut}},
  \tag{2d}
\end{equation}
using $h_{\min}=0.25$ (or the sharp geometric $h_{\min}$).  Combining
\eqref{eq:step2c},~\eqref{eq:step2d}, Lemma~\ref{lem:metric}, and
$\norm{\Delta e}_{L^2_{\mathrm{hyp}}}
\le A_{\mathrm{met}}\,C_{\Delta}\,\norm{e}_{H^2_{\mathrm{euc}}}$, and
writing $\norm{\lambda e}_{L^2}\le\lambda\norm{e}_{H^1}$, we obtain
\begin{equation}
\label{eq:step2e}
  \norm{(\Delta-\lambda)e_{\mathrm{bdry}}}_{L^2_{\mathrm{hyp}}}
  \le
  A_{\mathrm{ell}}\,
  \bigl(1+C_{\mathrm{Poincar\'e}}A_{\mathrm{met}}\bigr)\,
  \norm{e_{\mathrm{bdry}}}_{H^1_{\mathrm{hyp}}}
  =
  A_{\mathrm{res}}\,
  A_{\mathrm{bdry}}\,
  \delta_{\mathrm{aut}},
  \tag{2e}
\end{equation}
where $A_{\mathrm{ell}}=1+\lambda$ is a convenient majorant of the
elliptic/scaling constants on $K_Y$ (it absorbs both $\Delta e$ via the
$H^2\hookrightarrow H^1$ scale $h_{\min}^{-1}\le C$ and the zero-order
term $\lambda e$; equivalently $1+\lambda\ge C_{\Delta}(1+\lambda^{1/2})$
on $\lambda\ge 0$).  All $\lambda$-dependence of the corrector residual
lives here in $A_{\mathrm{res}}$, not in $A_{\mathrm{bdry}}$.  Therefore
\begin{equation}
\label{eq:step2}
  \norm{(\Delta-\lambda)u}_{L^2_{\mathrm{hyp}}(K_Y)}
  \le
  \tau_{\mathrm{tail}}
  +
  A_{\mathrm{bdry}}A_{\mathrm{res}}\,\delta_{\mathrm{aut}}
  \le
  A_{\mathrm{bdry}}A_{\mathrm{res}}\,\eta.
  \tag{2}
\end{equation}

\begin{remark}[No double count of $(1+\sqrt{\lambda})$]
\label{rem:no-double-sqrt}
The product $C_1\propto A_{\mathrm{bdry}}A_{\mathrm{res}}$ must not contain
both $(1+\sqrt{\lambda})$ and $(1+\lambda)$ for the same elliptic passage.
Step~1 controls only $\norm{e_{\mathrm{bdry}}}_{H^1}$ and is $\lambda$-free;
Step~2 lifts $H^1\to\norm{(\Delta-\lambda)e}$ and uses $A_{\mathrm{ell}}=1+\lambda$
once.  Dropping the spurious $(1+\sqrt{\lambda})$ from $A_{\mathrm{bdry}}$
improves $C_1$ by exactly the factor $1+\sqrt{\lambda}$ (at $r=6$:
$1+\sqrt{37}\approx 7.083$) without weakening any estimate in the proof.
\end{remark}

\paragraph{Step 3: collar / nonzero modes by Lemma~K.}
Extend $u$ into the collar $C_{Y_0}$ by the truncated Fourier expansion of
$\widetilde f$.  The omitted tail $N(\beta)>M$ contributes at most
$A_{\mathrm{cusp}}$ in $L^2_{\mathrm{hyp}}(C_{Y_0})$ by
Theorem~\ref{thm:lemmaK} and Lemma~\ref{lem:collar}.  The zero mode is
absent (cuspidality of the target subspace; Lax--Phillips~\cite{LP}).
A $C^2$ radial cutoff $\chi(y)$ equal to $1$ on $[Y,\infty)$ and $0$ on
$[Y_0/2,Y_0]$ (when $Y>Y_0$; else a fixed chart cutoff of width $Y_0/2$)
produces a periodizable function $U$ on $\Gamma\backslash\HH$ via
Shimizu's horoball lemma~\cite{Shimizu,Leutbecher}, with cutoff error
\begin{equation}
\label{eq:step3a}
  \norm{(\Delta-\lambda)\bigl((1-\chi)u\bigr)}_{L^2}
  \le
  A_{\mathrm{cut}}\,e^{-2\pi Y_0},
  \tag{3a}
\end{equation}
where $A_{\mathrm{cut}}$ as in~\eqref{eq:Acut} majorizes
$\abs{\chi''}+\abs{\chi'}$ times the $L^2$ mass of the first Fourier mode
in the collar (bounded by $\abs{T}Y^{-1}$ times Sobolev control of the
boundary value on $\{y=Y_0\}$).  Combining with~\eqref{eq:step2} and the
tail,
\begin{equation}
\label{eq:step3}
  \norm{(\Delta-\lambda)U}_{L^2(\Gamma\backslash\HH)}
  \le
  A_{\mathrm{bdry}}A_{\mathrm{res}}\,
  \bigl(1+A_{\mathrm{cusp}}\bigr)\,
  \eta
  +
  A_{\mathrm{cut}}\,e^{-2\pi Y_0}.
  \tag{3}
\end{equation}
(The factor $(1+A_{\mathrm{cusp}})$ accounts for replacing $u$ by its
$M$-truncated collar extension; when $A_{\mathrm{cusp}}\ll 1$ it is
essentially~$1$.  No extra $(1+C_K)$ appears: $C_K$ is already inside
$A_{\mathrm{cusp}}$.)

\paragraph{Step 4: spectral resolution and pigeonhole.}
Let $U_{\mathrm{cusp}}$ denote the orthogonal projection of $U$ onto the
cuspidal subspace (Step~5 controls $\norm{U-U_{\mathrm{cusp}}}$).  Write
\[
  U_{\mathrm{cusp}}
  =
  \sum_j c_j f_j
\]
in an $L^2$-orthonormal basis of cusp forms with eigenvalues
$\lambda_j=r_j^2+1$.  Self-adjointness gives
\begin{equation}
\label{eq:step4a}
  \sum_j \abs{\lambda_j-\lambda}^2\abs{c_j}^2
  =
  \norm{(\Delta-\lambda)U_{\mathrm{cusp}}}_2^2
  \le
  \norm{(\Delta-\lambda)U}_2^2.
  \tag{4a}
\end{equation}

\paragraph{Normalization under $\eta\le\eta_0$.}
From~\eqref{eq:step1} and~\eqref{eq:eta0},
$\norm{e_{\mathrm{bdry}}}_{H^1_{\mathrm{hyp}}}\le A_{\mathrm{bdry}}\eta
\le \tfrac14$ whenever $\eta\le\eta_0$.  With
$\norm{\widetilde f}_{L^2_{\mathrm{hyp}}(F_Y)}=1$ one gets
\begin{equation}
\label{eq:step4b}
  \norm{U}_{L^2}
  \ge
  \norm{u}_{L^2(F_Y)}
  -
  O\bigl(e^{-2\pi Y_0}\bigr)
  \ge
  1-\tfrac14-\tfrac18
  \ge
  \tfrac12
  \tag{4b}
\end{equation}
for $Y_0\ge 1/2$ and $\eta\le\eta_0$ (the $e^{-2\pi Y_0}$ collar loss is
$<\tfrac18$ at the engineering parameters; absorb into $\eta_0$ by
shrinking the first term of~\eqref{eq:eta0} if needed).  The cuspidal
projection loses at most the continuous mass controlled in Step~5, which
under $\eta\le\eta_0$ is $\le\tfrac18$; hence
$\norm{U_{\mathrm{cusp}}}\ge\tfrac12$ after possibly replacing $\eta_0$ by
the min in~\eqref{eq:eta0}.

\paragraph{Spectral pigeonhole.}
Set $R:=\norm{(\Delta-\lambda)U}_2$ and $\mu:=\norm{U_{\mathrm{cusp}}}_2\ge\tfrac12$.
From~\eqref{eq:step4a},
\begin{equation}
\label{eq:step4c}
  \sum_j \abs{\lambda_j-\lambda}^2\abs{c_j}^2
  \le
  R^2.
  \tag{4c}
\end{equation}
If every $\lambda_j$ satisfied $\abs{\lambda_j-\lambda}>R/\mu$, then
$\sum\abs{\lambda_j-\lambda}^2\abs{c_j}^2>(R/\mu)^2\sum\abs{c_j}^2=R^2$, a
contradiction.  Hence some $j_0$ obeys
\begin{equation}
\label{eq:step4d}
  \abs{\lambda_{j_0}-\lambda}
  \le
  \frac{R}{\mu}
  \le
  2R.
  \tag{4d}
\end{equation}
(Alternatively: the spectral measure of $U_{\mathrm{cusp}}$ has second
moment $\le R^2$ about $\lambda$, so its support intersects
$[\lambda-R/\mu,\lambda+R/\mu]$.)  Combining with~\eqref{eq:step3} and the
definition~\eqref{eq:C1}--\eqref{eq:C2},
\begin{equation}
\label{eq:step4}
  \abs{\lambda_{j_0}-\lambda}
  \le
  C_1\eta+C_2 e^{-2\pi Y_0},
  \tag{4}
\end{equation}
which is~\eqref{eq:lamerr} with $f:=f_{j_0}$ (up to a unit complex phase).

\paragraph{$L^2$ proximity.}
Let $P$ be the orthogonal projection onto $\mathrm{span}\{f_{j_0}\}$.  Then
\begin{equation}
\label{eq:step4e}
  \norm{U_{\mathrm{cusp}}-P U_{\mathrm{cusp}}}_2^2
  =
  \sum_{j\neq j_0}\abs{c_j}^2
  \le
  \frac{1}{\delta^2}
  \sum_{j\neq j_0}\abs{\lambda_j-\lambda}^2\abs{c_j}^2
  \le
  \frac{R^2}{\delta^2},
  \tag{4e}
\end{equation}
for any $\delta>0$ such that $\abs{\lambda_j-\lambda}\ge\delta$ for all
$j\neq j_0$.  Taking $\delta$ to be the isolation distance is unnecessary
for the crude bound: from~\eqref{eq:step4c} and
$\abs{c_{j_0}}^2=\mu^2-\sum_{j\neq j_0}\abs{c_j}^2$, the standard BSV
estimate yields
$\norm{U_{\mathrm{cusp}}-c_{j_0}f_{j_0}}_2\le 2R/\mu\le 4R$ under
$R\le\mu/2$ (which holds for $\eta\le\eta_0$ by~\eqref{eq:eta0} and
$C_1\eta\le\tfrac12\cdot\tfrac12$).  Absorbing constants into the factor
already present in $C_1$ gives~\eqref{eq:L2err}.

\paragraph{Step 5: continuous spectrum / cuspidal projection.}
Write $U=U_{\mathrm{cusp}}+U_{\mathrm{Eis}}$ via the Lax--Phillips
decomposition~\cite{LP}.  For $\lambda<1$ the continuous spectrum is
bounded away from $\lambda$ and
$\norm{U_{\mathrm{Eis}}}\le R/\mathrm{dist}(\lambda,[1,\infty))$.  For
$\lambda>1$ (Then regime $r\approx 6.6$) we state the theorem for the
cuspidal projection $U_{\mathrm{cusp}}$ as in~\eqref{eq:lamerr}--\eqref{eq:L2err}.
The projection error $\norm{U-U_{\mathrm{cusp}}}$ equals the Eisenstein
component; by the soft-cutoff argument of
BSV~\cite[Prop.~3.1]{BSV} and Child~\cite[Thm.~1.1]{Child}, specialized to
$\HH$ with the same collar cutoff as Step~3,
\begin{equation}
\label{eq:step5}
  \norm{U_{\mathrm{Eis}}}_2
  \le
  A_{\mathrm{Eis}}\,
  \bigl(
    \norm{(\Delta-\lambda)U}_2
    +
    \norm{U}_{H^{1/2}(\partial F_Y)}
  \bigr)
  \tag{5}
\end{equation}
with $A_{\mathrm{Eis}}\le A_{\mathrm{bdry}}$ (boundary traces of $U$ are
controlled by Step~1).  Thus the continuous mass is absorbed into the
same named constants; we do not claim a new soft-cutoff analysis, but
track the BSV/Child factor as $\le A_{\mathrm{bdry}}$ into $C_1$.

This completes the proof of Theorem~\ref{thm:DK}.
\endproof

\begin{remark}[What is not proved]
\label{rem:notproved}
Theorem~\ref{thm:DK} is strictly single-cusp.  It does not address:
\begin{itemize}[leftmargin=*]
\item two-cusp coupling or the block matrix
$\bigl(\begin{smallmatrix}V_\infty&-S\\-S^*&V_0\end{smallmatrix}\bigr)$;
\item congruence levels $N\mathfrak{p}$, reference cells, or gluing;
\item certified counting $N(\lambda)$, Weyl asymptotics as exact counts;
\item a uniform gap $\mu(N)$.
\end{itemize}
Once $C_1,C_2$ are fixed, the two-cusp extension is a block-matrix
estimate on top of the same defect machinery (Roadmap Rung~3).
\end{remark}

% ---------------------------------------------------------------------
\section{Numerical illustration}
\label{sec:numeric}

Then~\cite{Then} reports a first cuspidal spectral parameter for the
Picard group $\PSL_2(\ZZ[i])$ at
\[
  r_1 \approx 6.62212.
\]
As an \emph{engineering target} (not a theorem of this paper) we record
the interval
\begin{equation}
\label{eq:target}
  r_1\in[6.62211,\,6.62213].
\end{equation}
Converting~\eqref{eq:target} into a corollary of Theorem~\ref{thm:DK}
requires a concrete putative form $\widetilde f$ (e.g.\ from interval
Hejhal at $M=400$, $Y_0=0.8$) whose defect $\eta$ satisfies
\[
  C_1\eta+C_2 e^{-2\pi Y_0}
  \le
  10^{-5},
\]
so that $\lambda_{\mathrm{true}}$ is pinned inside the corresponding
$\lambda$-window.  The present paper supplies $C_1$ and $C_2$; the
production of $\widetilde f$ with certified $\eta$ is downstream work.

At the diagnostic parameters $(Y,Y_0,r)=(1.25,0.8,6.6)$ with sharp
$C_K=1$ and $y_{\min}=\tfrac12$, one has (cf.\ \texttt{constants\_DK.md})
\[
  C_1\sim 10^{5}\text{--}10^{6}
  \quad\text{(order; exact Arb values in the companion tables)},
\]
and the Lemma~K tail at $M=400$ is $<10^{-80}$, so the Fourier tail is
negligible compared to any realistic automorphy defect.  (Using the
classical $C_K^{\mathrm{cl}}(6.6)\approx 3.18\cdot 10^{4}$ inflates only
$A_{\mathrm{cusp}}$, which remains negligible at $M=400$.)

% ---------------------------------------------------------------------
\section{Bibliography}
\begin{thebibliography}{99}

\bibitem{BB11}
V.~Blomer and F.~Brumley,
\emph{On the Ramanujan conjecture over number fields},
Ann.\ of Math.\ (2) \textbf{174} (2011), 581--605.

\bibitem{BSV}
A.~R.~Booker, A.~Str\"ombergsson, and A.~Venkatesh,
\emph{Effective computation of Maass cusp forms},
Int.\ Math.\ Res.\ Not.\ 2006, Art.\ ID 71281.

\bibitem{Beb03}
M.~Bebendorf,
\emph{A note on the Poincar\'e inequality for convex domains},
Z.\ Anal.\ Anwendungen \textbf{22} (2003), 751--756.

\bibitem{Child}
S.~Child,
\emph{Verified computation of Maass cusp forms of arbitrary level and
eigenvalue},
PhD thesis / preprint (see also related computational spectral works).

\bibitem{DLMF}
NIST Digital Library of Mathematical Functions,
\url{https://dlmf.nist.gov/},
release~1.2.x; \S10.25--10.40 (modified Bessel functions).

\bibitem{EGM}
J.~Elstrodt, F.~Grunewald, and J.~Mennicke,
\emph{Groups Acting on Hyperbolic Space},
Springer Monographs in Mathematics, Springer, 1998.

\bibitem{Friedman}
J.~Friedman,
\emph{The Selberg trace formula and Selberg zeta-function for
cofinite Kleinian groups with finite-dimensional unitary representations},
Math.\ Z.\ (cf.\ related works on Bianchi spectral theory).

\bibitem{Iwaniec}
H.~Iwaniec,
\emph{Spectral Methods of Automorphic Forms},
2nd ed., AMS Graduate Studies in Mathematics \textbf{53}, 2002.

\bibitem{KS}
H.~H.~Kim (appendix by P.~Sarnak),
\emph{Functoriality for the exterior square of $\mathrm{GL}_4$ and the
symmetric fourth of $\mathrm{GL}_2$},
J.\ Amer.\ Math.\ Soc.\ \textbf{16} (2003), 139--183.

\bibitem{LP}
P.~D.~Lax and R.~S.~Phillips,
\emph{Scattering Theory for Automorphic Functions},
Ann.\ of Math.\ Studies \textbf{87}, Princeton Univ.\ Press, 1976.

\bibitem{Leutbecher}
A.~Leutbecher,
\emph{\"Uber die Heckeschen Gruppen~$G(\lambda)$},
Math.\ Ann.\ \textbf{146} (1962), 67--94.

\bibitem{Luke}
Y.~L.~Luke,
\emph{Inequalities for generalized hypergeometric functions},
J.\ Approx.\ Theory \textbf{5} (1972), 41--65.

\bibitem{PW60}
L.~E.~Payne and H.~F.~Weinberger,
\emph{An optimal Poincar\'e inequality for convex domains},
Arch.\ Rational Mech.\ Anal.\ \textbf{5} (1960), 286--292.

\bibitem{Shimizu}
H.~Shimizu,
\emph{On discontinuous groups operating on the product of the upper
half planes},
Ann.\ of Math.\ (2) \textbf{77} (1963), 33--71.

\bibitem{Then}
H.~Then,
\emph{Maass cusp forms for $\mathrm{PSL}(2,\mathbb{Z}[i])$},
Math.\ Comp.\ \textbf{77} (2008), 1701--1720? / computational tables.

\bibitem{Watson}
G.~N.~Watson,
\emph{A Treatise on the Theory of Bessel Functions},
2nd ed., Cambridge Univ.\ Press, 1944.

\bibitem{DualRoadmap}
Dual-certification roadmap (companion manuscript in this repository).

\end{thebibliography}

\end{document}
